\documentclass[11pt]{article}

\usepackage[margin=1in]{geometry}
\usepackage[T1]{fontenc}
\usepackage{amsmath,amssymb,amsthm}
\usepackage{array,longtable,booktabs}
\usepackage[hidelinks]{hyperref}

\newtheorem{definition}{Definition}
\newtheorem{lemma}{Lemma}
\newtheorem{theorem}{Theorem}
\newtheorem{proposition}{Proposition}
\newtheorem{corollary}{Corollary}

\newcommand{\Code}[1]{\texttt{#1}}
\newcommand{\Shape}{\mathsf{Shape}}
\newcommand{\Tele}{\mathsf{Tele}}
\newcommand{\Expr}{\mathsf{Expr}}
\newcommand{\Slot}{\mathsf{Slot}}
\newcommand{\Subst}{\mathsf{Subst}}
\newcommand{\Proper}{\mathsf{Proper}}
\newcommand{\nil}{\mathsf{nil}}
\newcommand{\pre}{\mathsf{pre}}
\newcommand{\dom}{\mathsf{dom}}
\newcommand{\cod}{\mathsf{cod}}
\newcommand{\mshape}{\mathsf{mid}}
\newcommand{\Act}[3]{#1^{\ast}_{#2}(#3)}
\newcommand{\Eta}{\eta}
\newcommand{\comp}{\mathbin{+\!\!+}}
\newcommand{\cons}{\mathbin{::}}
\newcommand{\ofList}{\mathsf{ofList}}
\newcommand{\one}{\langle - \rangle}

\title{Higher-Rank Syntax as a Relative Monad}
\author{Code companion for \Code{HigherRankSyntax/HigherRankSyntax/Equations.lean}}
\date{}

\begin{document}
\maketitle

\section{Purpose and influence}

This note explains the mathematical content of
\Code{HigherRankSyntax/HigherRankSyntax/Equations.lean}.  It is written in the
same spirit as the development of syntactic transformations in Anja Petkovic
Komel's thesis, especially the part beginning with the desiderata for
syntactic transformations and continuing through the relative-monad proof for
syntax.\footnote{\url{https://repozitorij.uni-lj.si/Dokument.php?id=152351&lang=slv}}

The guiding idea is the same.  First one defines a purely syntactic action on
expressions.  Then one proves that the action is compatible with the other
syntactic operations, especially instantiation/substitution.  Finally the
relative monad laws become short consequences of those interaction lemmas.

Our setting is a higher-rank generalization.  The role of a signature or
metacontext is played by a telescope-shaped object \(\Gamma : \Shape C\).  A
slot \(x : \Gamma \ni \alpha\) has an arity \(\alpha\), and each binder
\(j : C.\mathsf{Binder}\,\alpha\) carries its own arity \(j.\mathsf{arity}\).
Thus the application case is not simply a finite list of ordinary arguments:
each argument lives below a fresh binder specified by the head arity.  This is
the source of most of the proof engineering in \Code{Equations.lean}.

\section{Relative monads for this syntax}

The general relative-monad pattern is as follows.  The functor \(J\) gives the
generators of a syntactic context, while \(T\) gives expressions over that
context.  The unit embeds a generator as its generic expression.  A Kleisli map
assigns to each generator an expression in a target context, and its Kleisli
extension is the action of that assignment on all expressions.

\begin{definition}[Higher-rank generators and expressions]
For a shape \(\Gamma : \Shape C\), the generators of arity \(\alpha\) are slots
\[
  J\Gamma(\alpha) = \{x \mid x : \Gamma \ni \alpha\}.
\]
The corresponding expressions are
\[
  T\Gamma(\alpha) = \Expr(\Gamma \cons \alpha).
\]
The unit sends a slot to its eta-expansion:
\[
  \Eta_\Gamma(x) = \Code{Expr.eta}\;x.
\]
\end{definition}

\begin{definition}[Kleisli substitutions]
A Kleisli map from \(\Gamma\) to \(\Delta\) is a family
\[
  f : \prod_{\alpha}(\Gamma \ni \alpha) \to \Expr(\Delta \cons \alpha).
\]
In the Lean code this is wrapped as
\[
  \Code{toSubst}\;f : \Subst\,C\,\Shape.\nil\,\Gamma\,\Delta.
\]
The more general record
\[
  \sigma : \Subst\,C\,\pre\,\dom\,\cod
\]
keeps a fixed prefix \(\pre\), replaces slots from \(\dom\), and leaves the
prefix slots untouched:
\[
  \sigma(x) : \Expr(\pre \comp \cod \cons \alpha)
  \qquad(x : \dom \ni \alpha).
\]
\end{definition}

\begin{definition}[Substitution action]
For
\[
  \sigma : \Subst\,C\,\pre\,\dom\,\cod
\]
and a current binder depth \(\tau\), the action is
\[
  \Act{\sigma}{\tau}{-} :
  \Expr(\pre \comp \dom \comp \tau)
    \to
  \Expr(\pre \comp \cod \comp \tau).
\]
In Lean this is \Code{Subst.act}.  The depth \(\tau\) records binders already
entered by the recursive traversal.
\end{definition}

Mathematically, \Code{Subst.act} has three cases for an application head in
\(\pre \comp \dom \comp \tau\).

\begin{itemize}
\item If the head comes from \(\tau\), it is locally bound and is preserved.
\item If the head comes from \(\dom\), the substitution fires.
\item If the head comes from \(\pre\), it is outside the substitution domain
  and is only weakened through the new codomain.
\end{itemize}

The Lean classifier \Code{SubstSite} names exactly these cases:
\[
  \Code{tau},\quad \Code{dom},\quad \Code{pre}.
\]
The associated computation lemmas are
\Code{Subst.act\_ap\_inr}, \Code{Subst.act\_ap\_inl\_dom}, and
\Code{Subst.act\_ap\_inl\_pre}.

\section{Unit laws}

The first two relative-monad equations are the two unit laws.  In the thesis
style, these are proved after isolating the fact that the generic application
acts like the identity transformation.

\begin{lemma}[Action preserves eta at a bound head]
Let \(x : \tau \ni \alpha\).  Then
\[
  \Act{\sigma}{\tau \cons \alpha}
    {\Eta_{\pre \comp \dom \comp \tau}(x)}
  =
  \Eta_{\pre \comp \cod \comp \tau}(x).
\]
This is \Code{Subst.act\_eta\_inr}.
\end{lemma}

\begin{lemma}[Identity one-binder instantiation]
The identity instantiation for the singleton telescope \(\langle\alpha\rangle\)
acts as the identity on expressions.  This is the private recursion
\Code{idOf}, packaged as \Code{Subst.act\_inst.id}.
\end{lemma}

\begin{lemma}[Beta for eta]
Instantiating the eta-expansion of the singleton head by a one-binder kit
returns the corresponding kit application.  This is
\Code{Subst.act\_inst.eta}.
\end{lemma}

\begin{theorem}[Right unit]
For every expression \(e : \Expr(\Gamma \comp \tau)\),
\[
  \Act{\Code{Subst.id}\,\Gamma}{\tau}{e} = e.
\]
This is \Code{Subst.act\_id}.
\end{theorem}

\begin{theorem}[Left unit]
For a Kleisli map \(f\) and slot \(p : \Gamma \ni \alpha\),
\[
  \Act{\Code{toSubst}\,f}{\langle\alpha\rangle}{\Eta(p)} = f(p).
\]
This is \Code{Subst.act\_eta}.
\end{theorem}

\section{The main interaction lemma}

The thesis proves that a syntactic transformation commutes with
instantiation before using that fact to prove composition.  The direct analogue
in this file is the one-binder diamond:

\[
\begin{array}{ccc}
  \Expr((\pre\comp\dom\comp\tau\cons\alpha)\comp\ofList(\upsilon))
    & \xrightarrow{\Act{\sigma}{(\tau\cons\alpha)\comp\ofList(\upsilon)}{-}}
    & \Expr((\pre\comp\cod\comp\tau\cons\alpha)\comp\ofList(\upsilon))
\\
  \downarrow \kappa
    && \downarrow \kappa'
\\
  \Expr(\pre\comp\dom\comp(\tau\comp\ofList(\rho))\comp\ofList(\upsilon))
    & \xrightarrow{\Act{\sigma}{(\tau\comp\ofList(\rho))\comp\ofList(\upsilon)}{-}}
    & \Expr(\pre\comp\cod\comp(\tau\comp\ofList(\rho))\comp\ofList(\upsilon)).
\end{array}
\]

Here \(\kappa\) instantiates the source telescope with fillers
\(\iota\), while \(\kappa'\) instantiates it with the pointwise
\(\sigma\)-acted fillers.

\begin{theorem}[One-binder action/instantiation diamond]
Acting by \(\sigma\) and then instantiating the source telescope by
the acted fillers agrees with first instantiating that source telescope and then
acting by \(\sigma\).  In Lean:
\[
  \Code{Diamond.actThenInst}\;\sigma\;\rho\;\upsilon\;\iota\;e
  =
  \Code{Diamond.instThenAct}\;\sigma\;\rho\;\upsilon\;\iota\;e.
\]
This is \Code{diamondAt}.
\end{theorem}

The proof has five head cases, and these are the conceptual heart of the file.

\begin{center}
\begin{tabular}{lll}
\toprule
Case & Meaning & Proof behavior \\
\midrule
\Code{under} & the head is in \(\ofList(\upsilon)\) & preserve head, recurse on arguments \\
\Code{src} & the head is in the source telescope being instantiated & descend into the selected filler \\
\Code{tau} & the head is in the ambient depth \(\tau\) & preserve head, recurse on arguments \\
\Code{dom} & the head is substituted by \(\sigma\) & call the lifted companion \\
\Code{pre} & the head is in the fixed prefix & preserve/weaken head, recurse on arguments \\
\bottomrule
\end{tabular}
\end{center}

The \Code{dom} case is the reason the proof is mutual.  Substituting a
\(\dom\)-head exposes an expression \(\sigma(z)\), and the instantiation
around that expression must be rearranged.  This is handled by the companion
lemma.

\begin{theorem}[Lifted one-binder diamond]
Applying one singleton instantiation and then a second singleton instantiation
agrees with one fused singleton instantiation whose fillers already include
the first instantiation.  In Lean:
\[
  \Code{Lift.sequential}\;\rho\;\upsilon\;\chi\;\iota\;\eta\;e
  =
  \Code{Lift.fused}\;\rho\;\upsilon\;\chi\;\iota\;\eta\;e.
\]
This is \Code{liftAt}.
\end{theorem}

The companion has three head cases:
\[
  \Code{under},\qquad \Code{beta},\qquad \Code{pre}.
\]
The \Code{beta} case calls back into \Code{diamondAt}; ordinary argument
cases recurse structurally.

\section{Composition}

The substitution-level composition operation is
\[
  \Code{Subst.comp}\;\sigma\;\theta : \Subst\,C\,\pre\,\dom\,\cod,
\]
where
\[
  (\Code{Subst.comp}\;\sigma\;\theta)(x)
  =
  \Act{\theta}{\langle\alpha\rangle}{\sigma(x)}
  \qquad(x : \dom \ni \alpha).
\]
This is the general form of Kleisli composition, before specializing to empty
prefixes.

\begin{theorem}[Action preserves substitution composition]
For arbitrary substitutions
\[
  \sigma : \Subst\,C\,\pre\,\dom\,\mshape,
  \qquad
  \theta : \Subst\,C\,\pre\,\mshape\,\cod,
\]
and every expression \(e : \Expr(\pre\comp\dom\comp\tau)\),
\[
  \Act{\Code{Subst.comp}\;\sigma\;\theta}{\tau}{e}
  =
  \Act{\theta}{\tau}{\Act{\sigma}{\tau}{e}}.
\]
This is \Code{Subst.act\_comp}.
\end{theorem}

The proof follows the three head cases for \Code{Subst.act}.  The
\Code{tau} and \Code{pre} cases are preserved-head cases.  The \Code{dom}
case is the only substantive one: after \(\sigma\) substitutes the head, one
must commute \(\theta\) past the one-binder instantiation generated by the
arguments.  In the current file that bridge is a direct specialization of
\Code{diamondAt}, using the nil-target witness constructor
\Code{Proper.AppendCoherence.nil}.

\begin{corollary}[Kleisli composition law]
For Kleisli maps \(f : \Gamma \to T\Delta\) and
\(g : \Delta \to T\Xi\),
\[
  \Act{\Code{toSubst}(\Code{Subst.kcomp}\;f\;g)}{\tau}{e}
  =
  \Act{\Code{toSubst}\;g}{\tau}
    {\Act{\Code{toSubst}\;f}{\tau}{e}}.
\]
This is \Code{Subst.act\_kcomp}, now a short corollary of
\Code{Subst.act\_comp}.
\end{corollary}

\section{Code map}

\begin{longtable}{>{\raggedright\arraybackslash}p{0.27\textwidth}
                  >{\raggedright\arraybackslash}p{0.44\textwidth}
                  >{\raggedright\arraybackslash}p{0.21\textwidth}}
\toprule
Lean declaration & Mathematical role & Thesis-style analogue \\
\midrule
\Code{SubstSite} & Classifies an action head as \(\tau\), \(\dom\), or
\(\pre\). & Case split for syntax-directed action. \\

\Code{Subst.act\_ap\_inr} & Computation when the head is bound by the
current depth. & Variable/binder preservation case. \\

\Code{Subst.act\_ap\_inl\_dom} & Computation when substitution fires at a
domain head. & Nontrivial substitution case. \\

\Code{Subst.act\_ap\_inl\_pre} & Computation when the head is outside the
domain. & Parameter/prefix preservation case. \\

\Code{Subst.act\_inst.eta} & Beta law for eta-expanded singleton heads. &
Generic application computes correctly. \\

\Code{Subst.act\_inst.id} & Identity singleton instantiation acts as the
identity. & Identity instantiation is trivial. \\

\Code{Subst.act\_id} & Identity substitution acts as identity. & Relative
monad law \(\Eta^\dagger=\mathsf{id}\). \\

\Code{Subst.act\_eta} & Action of a Kleisli map on an eta-generator returns
the chosen filler. & Relative monad law \(f^\dagger\circ\Eta=f\). \\

\Code{diamondAt} & Action commutes with one-binder instantiation. &
Action interacts with instantiation. \\

\Code{liftAt} & Lifted companion needed by the domain branch of
\Code{diamondAt}. & Compatibility of nested instantiations. \\

\Code{Subst.comp} & Substitution-level Kleisli composition. & Composition of
syntactic transformations. \\

\Code{Subst.act\_comp} & Action preserves arbitrary substitution
composition. & Action preserves composition. \\

\Code{Subst.act\_kcomp} & Kleisli composition law for \Code{toSubst}. &
Third relative-monad law. \\

\Code{InterchangeFuel} & Well-founded measure for mutual one-binder
interaction. & Proof-technical structural induction measure. \\
\bottomrule
\end{longtable}

\section{What the Lean file should eventually say}

The current file now follows the intended mathematical ladder:

\begin{enumerate}
\item define the action and prove its three head computations;
\item prove the eta and identity-instantiation facts;
\item prove the one-binder interaction lemma, with its lifted companion;
\item prove arbitrary substitution composition;
\item specialize to the relative monad laws.
\end{enumerate}

The former \Code{PreNaturality} and \Code{Interchange} layers have been
deleted; their role is covered by the generalized \Code{diamondAt}/\Code{liftAt}
mutual theorem pair.

\end{document}
